\section{Rigid Fibre Drag Test Based on Marchetti}
\label{sec:marchetti-rigid-fibre-drag}

The settling-fibre comparison in Chapter~\ref{chap:coupled-validation} uses
the deformation data of \textcite{marchetti_deformation_2018} as a reference
for the coupled flexible-fibre model. The same paper also provides a useful,
more restricted, reference for a static rigid-fibre immersed-boundary test. The
purpose of such a test is not to reproduce the full flexible settling problem.
Instead, it isolates the hydrodynamic force computed by the immersed-boundary
scheme on a centerline representation of a slender body in a uniform flow.

The appropriate comparison is the weak-deformation limit discussed by
\textcite{marchetti_deformation_2018}. In that regime the fibre remains nearly
straight and horizontal while settling perpendicular to its long axis. The
hydrodynamic force is then approximated by the Stokes drag on a rigid slender
filament. Marchetti \emph{et al.} write the fibre length as \(2\ell\), the
radius as \(a\), and the aspect ratio as \(\kappa^{-1}=\ell/a\). With the
notation used here, \(L=2\ell\), \(D=2a\), and therefore
\(r_p=L/D=\ell/a=\kappa^{-1}\). The broadside rigid-fibre drag is
\begin{equation}
    F_{\perp}
    =
    C_{\perp}\mu U L,
    \qquad
    C_{\perp}
    =
    \frac{4\pi}{\log(4r_p)-1/2},
    \label{eq:marchetti-broadside-drag}
\end{equation}
up to the usual \(O(1/\log^2 r_p)\) slender-body correction. At leading order,
this reduces to \(C_{\perp}\simeq 4\pi/\log r_p\). Balancing
\eqref{eq:marchetti-broadside-drag} with the submerged weight gives the rigid
broadside settling speed \(U_{\perp}\) used by Marchetti \emph{et al.} to
nondimensionalize their measured velocities.

For the static-fibre test, the imposed flow speed is known rather than obtained
from a weight balance. The total force computed from the immersed-boundary
tractions can therefore be compared directly to
\eqref{eq:marchetti-broadside-drag}. If the drag coefficient is reported using
the projected area \(LD\),
\begin{equation}
    C_D
    =
    \frac{F_x}{\frac{1}{2}\rho U^2 LD},
    \qquad
    \Reynolds_D
    =
    \frac{\rho U D}{\mu},
\end{equation}
then the expected broadside target is
\begin{equation}
    C_{D,\perp}^{\rm SBT}
    =
    \frac{2C_{\perp}}{\Reynolds_D}
    =
    \frac{8\pi}
         {\Reynolds_D\left[\log(4r_p)-1/2\right]}.
    \label{eq:marchetti-broadside-cd-target}
\end{equation}
For the settling-fibre geometry used in this work, \(r_p=30\), giving
\(C_{\perp}=2.93\) and
\begin{equation}
    C_{D,\perp}^{\rm SBT}
    \simeq
    \frac{5.86}{\Reynolds_D}.
\end{equation}
This value is nearly identical to the finite-aspect-ratio perpendicular drag
formula used in the static-fibre template,
\begin{equation}
    C_{D,\perp}^{\rm TGT}
    =
    \frac{8\pi}
         {\Reynolds_D\left[\log(r_p)+\delta_{\perp}(r_p)\right]},
    \qquad
    \delta_{\perp}
    =
    0.839+\frac{0.185}{r_p}+\frac{0.233}{r_p^2},
\end{equation}
which gives \(C_D\simeq 5.92/\Reynolds_D\) for \(r_p=30\). The difference is
about one percent and is smaller than the expected discretization error for a
centerline immersed-boundary calculation.

This comparison should be made at a diameter Reynolds number small enough that
the Stokes drag law is the dominant physics. A practical sweep is therefore to
fix the same geometry and immersed-boundary spacing conventions as the settling
case,
\[
    r_p=30,
    \qquad
    \alpha_{\rm fibre}=\Delta s/h=2,
\]
and vary the finest immersed-boundary grid level. At each resolution, the
measured quantity is the time-averaged total hydrodynamic force on the pinned
rigid fibre after the uniform-flow solution has reached a statistically steady
state. The primary target is
\[
    \frac{C_D}{C_{D,\perp}^{\rm SBT}}\to 1
\]
as the fibre support width and marker spacing become small relative to the
diameter and length. Secondary diagnostics are the lift force, which should
remain small by symmetry, and the sensitivity of the drag to \(\Delta s/D\)
and the regularized-kernel support width relative to \(D\).

The figures and data in \textcite{marchetti_deformation_2018} provide context
for this test but not a direct immersed-boundary force distribution. Figure~2
shows the observed deformation regimes and identifies the small-\(B\) limit in
which the fibre remains close to a straight broadside filament. Figure~8 plots
the normalized maximum deflection against \(B\) and the elasto-viscous number
\(V\), and is the basis for the deformation data already digitized in
Figure~\ref{fig:marchetti-banaei-settling}. Figure~9 plots the measured and
bead-spring settling velocities as \(U/U_{\perp}\), giving the velocity target
\(U/U_{\perp}\simeq 1\) in the weak-deformation rigid limit. Figure~10 plots
the dimensionless drag relation between \(B\) and \(V\), showing a
weak-deformation branch where \(B\simeq C_{\perp}V\), an intermediate
reconfiguration branch where the drag is shape dependent, and a saturated
branch approaching the drag of two nearly vertical filament arms. These figures
justify the broadside rigid limit as a force target, but they do not provide
local hydrodynamic tractions or a finite-Reynolds-number drag law for a
centerline fibre.

The static test should therefore be interpreted as a low-Reynolds-number
hydrodynamic-force verification of the fibre immersed-boundary representation.
It is complementary to the full settling validation at
\(\mathrm{Ga}=40\). With the Banaei nondimensionalization used for the settling
case and \(r_p=30\), the velocity scale corresponds to a diameter Reynolds
number of approximately \(\Reynolds_D=\mathrm{Ga}/r_p\simeq 1.33\). That value
is low, but it is not asymptotically Stokesian. A static run at
\(\Reynolds_D\simeq 1.33\) can be useful as a Ga-comparable sensitivity check,
whereas the strict Marchetti slender-body target should be assessed using
\(\Reynolds_D\ll1\), for example \(\Reynolds_D=0.05\) or \(0.1\).
